\chapter{往期错题}

\begin{mistake}
计算：$\dfrac{4}{7}\times\dfrac{14}{15}\div\dfrac{3}{5}=\dfrac{(\quad)}{(\quad)}\times\dfrac{(\quad)}{(\quad)}\times\dfrac{(\quad)}{(\quad)}=\dfrac{(\quad)}{(\quad)}=\dfrac{(\quad)}{(\quad)}$
\end{mistake}

\vspace{5cm}
\begin{mistake}
$\dfrac{7}{13}\times\dfrac{5}{14}\times\dfrac{39}{8}$
\end{mistake}

\vspace{5cm}
\begin{mistake}
$\dfrac{3}{5}\div6\times\dfrac{4}{15}$
\end{mistake}

\vspace{5cm}
\begin{mistake}
$\dfrac{11}{26}\div\dfrac{5}{13}$
\end{mistake}

\vspace{5cm}
\begin{mistake}
小红在计算$7\times(\blacksquare+0.05)$时错算成$7\times\blacksquare+0.05$，这样算出的结果与正确的结果相差（$\quad$）
\end{mistake}

\vspace{5cm}
\begin{mistake}
一个房间放了两张桌子，第一张桌子是长方形的，长1.4米，宽0.6米，第二张桌子是正方形的，边长是0.8米，哪张桌子的面积大一些，大多少平方米？
\end{mistake}

\vspace{5cm}
\begin{mistake}
李叔叔需要若干根长0.36米的短绳，他找来一根长绳，截了八次，正好截完，这根绳子原来长多少米？
\end{mistake}

\vspace{5cm}
\begin{mistake}
把五个数按照从小到大的顺序排成一列，前三个数的平均值是0.52，后三个数的平均值是0.86，这五个数的平均数是0.65，这列数最中间的数是多少？
\end{mistake}

\vspace{5cm}
\begin{mistake}
$4.1\times7.6=$
\end{mistake}

\vspace{5cm}
\begin{mistake}
得数保留两位小数：$5.89\times3.6\approx$
\end{mistake}

\vspace{5cm}
\begin{mistake}
用13.2米的篱笆一面靠墙围成一块宽3.2米的长方形菜地，这块地的面积最大是多少平方米？
\end{mistake}

\vspace{5cm}
\begin{mistake}
$6.98\times4.06$的积与下面算式中（$\quad$）的积最接近

A. $6\times5$\quad B. $7\times4$\quad C. $7\times5$\quad D. $6\times4$
\end{mistake}

\vspace{5cm}
\begin{mistake}
六（1）班42名同学合影留念，全班每人洗一张照片，一共需要花多少钱？（原稿缺价格信息，待补）
\end{mistake}

\vspace{5cm}
\begin{mistake}
妈妈带40元去市场买鸡，每千克鸡19.5元。鸡重1.9千克，妈妈带的钱够吗？
\end{mistake}

\vspace{5cm}
\begin{mistake}
一列火车长160米，以每分钟0.95千米的速度通过一座大桥，从车头上桥到最后一节车厢离开桥共用了2.6分钟，这座大桥长多少米？
\end{mistake}

\vspace{5cm}
\begin{mistake}
李奶奶家上半年用水34.2吨，王奶奶家一季度一共用水25.8吨。谁家用水多？多多少？
\end{mistake}

\vspace{5cm}
\begin{mistake}
明明和军军拿出同样的钱合买一箱苹果，明明分得5.8千克，军军分得4.2千克。最后明明给了军军8.4元。问苹果多少钱每千克？
\end{mistake}

\vspace{5cm}
\begin{mistake}
$38.0202\ldots$精确到十分位，结果是（$\quad$）

A. 38\quad B. 38.0\quad C. 38.1\quad D. 38.2
\end{mistake}

\vspace{5cm}
\begin{mistake}
在$\bigcirc$里填“$>$”“$<$”或“$=$”。

$0.54\div1.2\bigcirc0.54\div0.75$
\end{mistake}

\vspace{5cm}
\begin{mistake}
先计算，再按要求写出商的近似值：

\begin{center}
\begin{tabular}{|c|c|c|c|}
\hline
 & 保留一位小数 & 保留两位小数 & 保留三位小数\\
\hline
$6.5\div1.6$ & & & \\
\hline
$27\div11$ & & & \\
\hline
$4.4\div0.23$ & & & \\
\hline
\end{tabular}
\end{center}
\end{mistake}

\vspace{5cm}
\begin{mistake}
一种纸箱最多可以装35个鸡蛋，用这种纸箱装2422个鸡蛋，至少需要准备多少个纸箱？
\end{mistake}

\vspace{5cm}
\begin{mistake}
六（1）班教室宽7米，长是宽的1.4倍。

（1）如果六（1）班有学生42人，教室的人均使用面积大约是多少平方米？（得数保留两位小数）

（2）如果按照《中小学校设计规范》的规定，小学普通教室的人均使用面积应大于或等于1.36平方米。这间教室最多可以容纳多少学生？
\end{mistake}

\vspace{5cm}
\begin{mistake}
十寸为一尺，十尺为一丈，三国时，一尺约为0.24米，张飞使用的“丈八蛇矛”长一丈八寸，约（$\quad$）米长。
\end{mistake}

\vspace{5cm}
\begin{mistake}
把10枚1元硬币整齐地叠放在一起，高度大约是1.85厘米，照这样计算，1000枚一元硬币叠放在一起的高度大约是（$\quad$）厘米，一亿枚一元硬币叠放在一起大约是（$\quad$）千米。
\end{mistake}

\vspace{5cm}
\begin{mistake}
算式$3\square.08\times\square0.8$的积可能是（$\quad$）

A. 615.24\quad B. 9.64\quad C. 63.44\quad D. 610.524
\end{mistake}

\vspace{5cm}
\begin{mistake}
一盒果汁有1升，先倒满3小瓶，每小瓶$\dfrac{1}{4}$升，剩下的平均倒入2个杯子，每杯多少升？
\end{mistake}

\vspace{5cm}
\begin{mistake}
$6.4\div0.16$\quad $7.8\div2.3$（保留两位小数）
\end{mistake}

\vspace{5cm}
\begin{mistake}
$2.03\times0.46$（得数保留两位小数）
\end{mistake}

\vspace{5cm}
\begin{mistake}
$(1.2+1.3)\div0.25$
\end{mistake}

\vspace{5cm}
\begin{mistake}
$\dfrac{1}{2}+\dfrac{5}{8}\square4$（运算符号原稿缺失，待补）
\end{mistake}

\vspace{5cm}
\begin{mistake}
$11.57-0.4+8.43-9.6$
\end{mistake}

\vspace{5cm}
\begin{mistake}
一种葡萄的单价是7.8元/千克。小华和小丽各带了35元，小华想买5斤葡萄，小丽想买3斤葡萄。如果各自付钱，两人带的钱都够吗？如果合起来付钱，一共要付多少元？
\end{mistake}

\vspace{5cm}
\begin{mistake}
$4.26\times0.31$的积是（$\quad$），$7.8\div0.12$的商的最高位在（$\quad$）位；$15\div\square$的商用循环小数表示为（$\quad$）。（除数原稿缺失，待补）
\end{mistake}

\vspace{5cm}
\begin{mistake}
珂珂计算一个算式时，把一个数除以4.2算成了乘4.2，得到的结果是26.46，这个算式的正确结果是（$\quad$）
\end{mistake}

\vspace{5cm}
\begin{mistake}
趣味运动会上，队员需要走完96米的赛道，平均每一步迈出0.75米，走完全程一共用时68秒。走完这段赛道，队员迈出一步需要（$\quad$）秒。（保留一位小数）
\end{mistake}

\vspace{5cm}
\begin{mistake}
2026年冬奥会和冬残奥会的吉祥物是娜咖和米罗。一个两层货架，上层放的吉祥物玩偶数量是下层的2.5倍，如果从上层取60个玩偶放入下层，那么两层吉祥物玩偶的个数相等。上层原来有（$\quad$）个玩偶。
\end{mistake}

\vspace{5cm}
\begin{mistake}
乐乐列竖式计算“$7.76\div2.5$”的商。如图，当商“3.1”余数为“1”时，这里的“1”表示（$\quad$）（原题配有余数竖式图）

A. 1个一\quad B. 1个十分之一\quad C. 1个百分之一\quad D. 1个千分之一
\end{mistake}

\vspace{5cm}
\begin{mistake}
一列火车长180米，以每分钟0.9千米的速度通过一座大桥，从车头上桥到最后一节车厢离开大桥共用了1.8分钟，这座大桥长多少千米？
\end{mistake}

\vspace{5cm}
\begin{mistake}
薯片0.8元/袋，用9.6元可以买几袋？列式为（$\quad$）$\div$（$\quad$）$=$（$\quad$）
\end{mistake}

\vspace{5cm}
\begin{mistake}
学校里有个块长3.6米、宽2.4米的长方形宣传栏（如图1）。依次通过列竖式（如图2）计算它的面积。①竖式“$4\times6$”表示4个（$\quad$）和6个（$\quad$）相乘，所得的积表示宣传栏中（$\quad$）（填序号）的面积。“72”表示72个（$\quad$），指宣传栏中（$\quad$）（填序号）的面积是（$\quad$）平方米。（原题配图1、图2）
\end{mistake}

\vspace{5cm}
\begin{mistake}
$2.56\times9.2=$\quad $0.75\div0.18\approx$（保留两位小数）
\end{mistake}

\vspace{5cm}
\begin{mistake}
泉州美食价格如下：面线糊18.5元/碗，满煎糕9.5元/斤，礼盒12.5元/个（每盒装2.5斤）。

（1）80元最多买多少碗面线糊？（2）买8斤满煎糕都用礼盒装，买满煎糕和礼盒最少花多少钱？
\end{mistake}

\vspace{5cm}
\begin{mistake}
王阿姨开小汽车上班，每行驶1km大约耗油0.07升。她家距离单位15km，每月按22天计算，每天往返一次，每月耗油多少升？油价6.94元/升，每月的油费大约是多少元？（保留两位小数）
\end{mistake}

\vspace{5cm}
\begin{mistake}
开学果园采摘了0.8吨梨，分在大、小两种袋子中售卖（小袋8.4元/0.8kg，大袋13.8元/1.5kg）。

（1）哪种规格的梨每千克售价更低？（2）如果把今天的梨都装进大袋子中，至少需要多少个大袋子？
\end{mistake}

\vspace{5cm}
\begin{mistake}
$4.5-\dfrac{6}{7}=$
\end{mistake}

\vspace{5cm}
\begin{mistake}
$12.5\times3.2\times0.25=$\quad $15\times\left(\dfrac{4}{15}+\dfrac{3}{19}\right)\times19=$
\end{mistake}

\vspace{5cm}
\begin{mistake}
有一桶油，油重15千克，第一次倒出它的$\dfrac{2}{5}$，第二次倒出$\dfrac{2}{5}$千克，剩下的油正好可以倒满4个瓶子，平均每个瓶子可以装多少千克油？
\end{mistake}

\vspace{5cm}
\begin{mistake}
某工程队10天挖好一条水渠。前3天平均每天挖一千米，后7天一共挖了7.2千米，这个工程队平均每天挖水渠多少千米？
\end{mistake}

\vspace{5cm}
\begin{mistake}
看图列式计算（图：六（1）班48人；六（2）班比六（1）班少$\dfrac{1}{6}$，求六（2）班多少人）
\end{mistake}

\vspace{5cm}
\begin{mistake}
商场购进一批电视机，先提价$\dfrac{1}{10}$后降价$\dfrac{1}{10}$销售，现价（$\quad$）原价。

A. 大于\quad B. 小于\quad C. 等于
\end{mistake}

\vspace{5cm}
\begin{mistake}
12米增加它的$\dfrac{1}{3}$后，再减少$\dfrac{1}{3}$米，是（$\quad$）米。
\end{mistake}

\chapter{中秋练习}
\newpage
\begin{mistake}
$0.68\times 0.49$的积是(  )位小数\\
$8.5\div 0.06$转化成整数除法是(   )\\商的最商位是(  )位。
\end{mistake}

\vspace{5cm}
\begin{mistake}
小轩在计算 $2.3\times 1.2$时，用的是下面图中的方法:

\begin{center}
\begin{tikzpicture}[font=\normalsize]
    % 左：小轩的计算过程
    \draw[thick] (0,1.35) rectangle (3.6,-1.35);
    \node[anchor=north west, align=left, inner sep=0pt,
          font=\normalsize\linespread{1.1}\selectfont] at (0.32,1.12)
        {$2.3\times1.2$\\[2pt]
         $=2\times1+0.3\times0.2$\\[2pt]
         $=2+0.06$\\[2pt]
         $=2.06$};

    % 右：面积模型
    \begin{scope}[shift={(5.3,1.35)}]
        \draw[thick] (0,-2.7) grid[xstep=1.7,ystep=0.9] (3.4,0);
        \node at (0.85,-0.45) {$2$};
        \node at (2.55,-0.45) {$0.3$};
        \node at (-0.5,-1.35) {$1$};
        \node at (-0.5,-2.25) {$0.2$};
        \node at (0.85,-1.35) {$2\times(\quad)$};
        \node[LiSiWarn] at (2.55,-1.35) {$(\quad)\times(\quad)$};
        \node[LiSiWarn] at (0.85,-2.25) {$(\quad)\times(\quad)$};
        \node at (2.55,-2.25) {$0.3\times0.2$};
    \end{scope}
\end{tikzpicture}
\end{center}
1. 请根据算式$2.3\times 1.2$，在右图中填一填。

2.小轩少算了哪些部分?($\quad$)(填写算式)

3.正确结果是($\quad$)
\end{mistake}
\vspace{8cm}
\begin{mistake}
辆越野车在沙漠中行驶了 $32.5$ 千米，耗油 $5.2$ 升。如果这辆越野车要穿越总路程为$1309$ 千米的沙漠无人区，至少得准备(   )升汽油。(得数保留整数)
\end{mistake}

\vspace{5cm}
\begin{mistake}
$m\times n=\dfrac{1}{200}$，$(m\times100)\times(n\div20)=$（$\quad$）；$m+n=150$，$m\div\dfrac{5}{3}+n\times0.6=$（$\quad$）。
\end{mistake}

\vspace{5cm}
\begin{mistake}
乐乐家书架上的书的本数在 $40\sim50$ 本之间，其中科普书占$\dfrac{1}{6}$，故事书占$\dfrac{3}{8}$，其余的是教辅书，教辅书有（$\quad$）本。
\end{mistake}

\vspace{5cm}
\begin{mistake}
列竖式计算：$2.56\times 9.2$
\end{mistake}

\vspace{5cm}
\begin{mistake}
计算：$0.42\div[(3.2-0.5)\div 9]$
\end{mistake}

\vspace{5cm}
\begin{mistake}
儿童负重最好不要超过体重的$\dfrac{3}{20}$，小贝的体重是35千克，书包重5千克，小贝的书包（$\quad$）

A. 超重了\quad B. 没有超重\quad C. 刚好符合\quad D. 无法判断
\end{mistake}

\vspace{5cm}
\begin{mistake}
猎豹平均每小时可跑80千米，按照这个速度，一只猎豹从草原上的A点跑到B点一共需要2小时，$\underline{\qquad\qquad}$，它第2个小时跑了多少千米？列式为$80\times2\times\left(1-\dfrac{3}{5}\right)$，那么横线上应该补充的条件是（$\quad$）

A. 第2个小时跑的距离比第1个小时少$\dfrac{3}{5}$\\
B. 第1个小时跑的距离比第2个小时少$\dfrac{3}{5}$\\
C. 第1个小时跑了全程的$\dfrac{3}{5}$\\
D. 第2个小时跑了全程的$\dfrac{3}{5}$
\end{mistake}

\vspace{5cm}
\begin{mistake}
一辆大巴前往红色教育基地，去时每小时行56千米，$\dfrac{9}{8}$小时到达；原路返回时只用了0.375小时。这辆大巴往返一次的平均速度是多少千米/时？
\end{mistake}

\vspace{5cm}
\begin{mistake}
加工一批零件，师傅和徒弟合作6天可以完成。现在师傅先加工4天，徒弟接着加工6天，完成了这批零件的$\dfrac{4}{5}$。照这样计算，徒弟单独加工这批零件需要多少天？
\end{mistake}

